\section{Policy objective, comparators, and heterogeneity}
\label{sec:policy}

The experimental sham/fault sampling ratio does not define deployment value.
For deployment fault prevalence $p$,
\begin{equation}
 V(d;p)=pV_1(d)+(1-p)V_0(d),\qquad
 G(d;p)=pG_1(d)+(1-p)G_0(d).
 \label{eq:prevalence-value}
\end{equation}
Cost and invalid-output curves use the same mixture.  Curves will be reported
for $p\in[0,1]$, with $p\in[0,0.25]$ marked only as the provisional plausible
operating region.  Values outside that interval are possible, not forbidden.

Candidate policies come from a finite, prospectively specified class $\Dcal$
mapping $(Z,M)$ to an available regime; a policy is eligible only if held-out
estimates certify \cref{eq:noninferiority,eq:cost-constraint,eq:invalid-constraint}
simultaneously (\cref{prop:certification}).  The learned policy $d^\dagger$ is
the gated deviation from the non-personalized comparator of
\eqref{eq:gated-rule}, and \cref{thm:safe} gives the sense in which it cannot do
worse.  This is finite-class held-out policy selection under availability
masking, not a new safe reinforcement-learning algorithm.

The comparator hierarchy is: $g_0$; the best fixed feasible intervention rule;
a random matched-cost intervention; failure-risk-gated verification;
effect-gated verification; an \emph{optimistic} controller using point estimates
of regime value from the selection split, greedily, without the lower confidence
bound; the two parametric policy candidates, the \emph{unconstrained policy}
and \emph{constrained policy}; and the two derived deployment rules, the
\emph{certified constrained policy} (the constrained policy under
the pessimistic gate of \cref{cor:proposal}) and the learned multi-action policy
$d^\dagger$ of \eqref{eq:gated-rule}.  The optimistic arm isolates the
contribution of pessimism: if optimistic and gated controllers perform
identically, the lower confidence bound is doing no work and \cref{thm:safe} is
decoration; if the optimistic controller wins on mean completion while breaching
the no-fault tolerance, that is the trade-off the theorem exists to manage,
reported in those terms.  Every comparator uses the same availability mask,
prevalence, budgets, constraints, and final data.
Because a single named regime may be unavailable at some checkpoints, the
strongest non-personalized comparator is an availability-aware fixed-priority
rule selected from a registered class $\Dcal_{\mathrm{fixed}}$:
\begin{equation}
 d^\star_{\mathrm{fixed}}(p)
 =\argmax_{d\in\Dcal_{\mathrm{fixed}}\,:\,d\text{ certified}} V(d;p).
\end{equation}
The $\argmax$ is evaluated on the calibration split, so the comparator is frozen
before final evaluation.  It is $p$-indexed: \eqref{eq:personalization}
therefore compares the gated policy against a \emph{different} fixed rule at
each prevalence, which is the honest reading of ``the strongest fixed
rule at this deployment condition'' rather than a single winner chosen once.
The confirmatory personalization estimand is
\begin{equation}
 G_{\mathrm{personal}}(p)
 =V(d^\dagger;p)-V(d^\star_{\mathrm{fixed}};p).
 \label{eq:personalization}
\end{equation}

\paragraph{The registered candidate class.}
The classes are finite and are enumerated here; nothing else enters the
multiplicity count of \cref{prop:certification}.  The non-personalized class
$\Dcal_{\mathrm{fixed}}$ contains six fixed, availability-aware rules: $g_0$;
always-verify; always-recover (the family-specific recovery regime when
available); a random matched-cost intervention; verification gated on a
predicted baseline failure risk above a registered threshold; and verification
gated on an estimated verification effect above a registered threshold.  The
candidate class $\Dcal$ is $\Dcal_{\mathrm{fixed}}$ together with three
policies whose realized form is determined before the calibration data are
observed: the optimistic controller (its point estimates of regime value come
from the selection split), and the two parametric policies fit by the one-step
learner of \cref{sec:rl}, the unconstrained and constrained policy members.
Architecture, hyperparameters and checkpoint are fixed on the selection split.
So $|\Dcal|=9$.  Each learned policy is one frozen policy; internal hyperparameter
and checkpoint variants are selected within the selection split and do not
inflate the class.

Two deployment rules are \emph{derived}, not members of $\Dcal$: $d^\dagger$
and the certified constrained policy, the constrained policy with
the pessimistic gate of \cref{cor:proposal} applied on calibration data.  Their
realized form depends on the calibration bounds, so they carry no
pre-registered value bound; their safety is inherited from
\cref{thm:safe,cor:proposal} rather than from membership, their constraints are
measured directly, and their value is reported from the untouched final split as
\cref{sec:stages} prescribes.  In the terms of \cref{as:frozen}, members of
$\Dcal$ are realized policies determined before the calibration data are
observed, and a rule obtained by applying a registered gate to a proposal is an
output of the procedure rather than a member.  \emph{Constrained} is not
\emph{certified}: the constraints of the constrained policy hold
on the training split only, and certification is the separate step of
\cref{prop:certification} that applies to whichever member of $\Dcal$ is
eventually proposed.  The oracle branch-selection upper bound is an evaluation
reference and is not a member of $\Dcal$.

\subsection{One-step policy learning}
\label{sec:rl}

The learned policy is the strongest \emph{opponent} the certification of
\cref{sec:experiments} is measured against, not the method of this paper: it
supplies none of the mathematics --- that is
\cref{thm:ordinal,thm:refinement,thm:safe,prop:certification} --- and it is
precisely the thing that would break the no-fault line, which is why
$\Delta_{\mathrm{certification}}$ of \cref{sec:experiments} is worth measuring.
Exhaustive branches give action-specific outcome, harm, invalidity and cost
feedback for every feasible arm, which is what makes such an opponent trainable
at all.

Branching makes the intervention value of every feasible action quantitative at
each checkpoint, so the opponent is trained by \emph{masked full-information
contextual policy learning}.  This is a one-step structured decision with
action-specific feedback for every feasible registered regime, not a
token-level relative-policy objective.  The context at checkpoint $i$ is
$x_i=(Z_i,M_i)$; the action is a regime $g\in\Gset(H_i)$; the frozen
continuation and the branch design supply terminal outcomes $Y^g_i$, $I^g_i$ and
$C^g_i$ for every feasible $g$.  A masked policy over the structured decision
space is
\begin{equation}
 \pi_\theta(g\mid Z,M)\;\propto\;M_{g}\,\exp s_\theta(Z,g),
 \label{eq:masked-policy}
\end{equation}
with the mask applied to the logits, where $s_\theta$ is a small MLP or linear
score over the registered representation --- not a token-level language model,
because the decision space is a handful of structured actions.

The training signal is the checkpoint-relative advantage.  Let
$\pi_{\mathrm{ref}}$ be a registered reference distribution over the feasible
set (uniform, or the comparator-induced distribution), and let
$\widehat\mu_{ig}$ be the cell-level estimate of $\mu^{(J_i)}_g$ at checkpoint
$i$.  The advantage of regime $g$ at $i$ is
\begin{equation}
 \widehat A_{ig}=\widehat\mu_{ig}-\sum_{h}\pi_{\mathrm{ref}}(h\mid x_i)\,\widehat\mu_{ih},
 \label{eq:checkpoint-advantage}
\end{equation}
Because $\sum_g\pi_{\mathrm{ref}}(g\mid x_i)=1$, these advantages are zero-sum
within a checkpoint,
$\sum_g\pi_{\mathrm{ref}}(g\mid x_i)\widehat A_{ig}=0$: the horizon-one,
regime-level analogue of the step-shared zero-sum shaping of
\citet[Prop.~2]{adrs2026}.\footnote{No potential-based invariance argument is
needed here. At horizon one there is no downstream policy for a potential to be
invariant under, because the frozen continuation is part of every regime and the
branch is the only backup.}
The two terms are \emph{cross-fitted}: repetitions within each
checkpoint--regime cell are split, the selected term estimated on one half and
the evaluated term on the other.  Using the same $\widehat\mu$ in both terms is
selection on noise --- the identical mechanism that produced a large positive bias
under the null in \cref{sec:experiments}.  The objective is
\begin{equation}
 \widehat J(\theta)=\frac1N\sum_i\sum_{g}\pi_\theta(g\mid Z_i,M_i)\,\widehat A_{ig},
 \label{eq:rl-objective}
\end{equation}
summed over the feasible set only.  Full feedback holds \emph{conditional on
pre-treatment feasibility}, $\Prb(E_{ig}=1\mid A_{ig}=1)=1$; masked arms are
restricted, not missing, so no action-assignment positivity must be defended on
the feasible set.

It is worth being exact about what \eqref{eq:rl-objective} is, because the
Bellman reading of \cref{app:notation} invites a wrong identification.  The
objective is a \emph{policy score}: $\theta$ enters only through the masked
softmax \eqref{eq:masked-policy}, and the advantages $\widehat A_{ig}$ are
plug-in cell-level quantities held fixed during the optimisation.  The
procedure therefore never regresses $Q^\rho$ onto its inputs, never forms a
temporal-difference residual, and is not an action-value network.  A fitted
$Q^\rho$ regression --- one model trained on full-arm terminal outcomes with
the controller's observation $(z,m)$ and the regime as input --- is a coherent
alternative that would need no bootstrap term either, since at horizon one the
branch \emph{is} the backup; it is proposed in \cref{app:analysis} as an
additional baseline rather than as a description of this learner, it is not a
member of $\Dcal$, and a variant given the full state $X$ would be an oracle
diagnostic rather than a deployable comparator.  The cell-level
estimates $\widehat\mu_{ig}$ that feed \eqref{eq:checkpoint-advantage} are
produced by the frozen score library --- penalised logistic regression and
shallow gradient-boosted trees, hyperparameters chosen from a frozen grid by
template-grouped inner cross-validation --- which fits a probability for the
risk score and a real-valued paired advantage for the effect score.  That
library is what the repository implements; nothing in it is a $Q$-learner.  Although $J_i$ is not a deployment feature, it enters the
objective through the stratum mixture used to evaluate policies at a prevalence
point, $\widehat V(\pi;p)=p\widehat V_1(\pi)+(1-p)\widehat V_0(\pi)$.  The
procedure is a standard masked full-information contextual policy learner; the novelty is the
branch design and the constraints, not the optimizer.  The policy is fit on the
training split, hyperparameters and checkpoint fixed on the selection split, and
the result is the unconstrained-policy member of $\Dcal$.

The constrained candidate maximizes the same value subject to the no-fault,
invalid-output and vector-cost constraints of \cref{sec:constraints}, via the
Lagrangian
\begin{equation}
 \max_\theta\ \widehat V(\pi_\theta;p)-\lambda_0\,\hat g(\pi_\theta)-\lambda_1\,\hat r(\pi_\theta)-\sum_k\lambda_{2k}\,\hat c_k(\pi_\theta),
 \label{eq:rl-lagrangian}
\end{equation}
where $\hat g,\hat r,\hat c_k$ are the training-split violations of
\cref{eq:noninferiority,eq:invalid-constraint,eq:cost-constraint}.  The
multipliers are hyperparameters fixed on the \emph{selection split}: tuning them
against calibration data would make the constrained policy calibration-dependent
and violate \cref{as:frozen}.  \emph{Constrained} is not \emph{certified}: the
constraints hold on the training split only, and certification is the separate
held-out step of \cref{prop:certification}.  The result is the
constrained-policy member of $\Dcal$.

\subsection*{The actionability gap (deferred)}

\emph{Status.} The two studies named here are deferred and this subsection is a
retained design description, not a result.  Supporting study E2 would
characterize how far baseline failure risk and actionable value
\eqref{eq:actionable} come apart; study E3 would test whether the
separation makes risk ordinally insufficient for allocating verification.  Three
regimes are distinguishable, and E2 would report their mass by fault family,
propagation depth, remaining budget, availability, and verifier reliability:

\begin{enumerate}[label=(\roman*),leftmargin=2.2em]
\item \emph{High baseline risk, low actionable value.}  The trajectory is
  already beyond the reach of any available regime, so a warning is correct and
  useless.
\item \emph{Moderate baseline risk, high actionable value.}  A cheap available
  intervention prevents propagation, so the most valuable warning is one a
  risk-ranked monitor would suppress.
\item \emph{Low baseline risk, individually harmful regimes.}  The checkpoint is
  clean and some available regime damages it; this is
  \eqref{eq:harm-fraction} restricted to $J=0$.
\end{enumerate}

Regime (iii) is invisible to any procedure whose output is a failure
probability, because such a procedure has no representation for the damage an
intervention does.  The external evidence for this heterogeneity is given in
\cref{sec:main-introduction}: it motivates, but does not identify, these
checkpoint-level regimes.  The prediction is that all three regimes carry
non-negligible mass and that regime (iii) is non-empty in the no-fault stratum.
The deferred gap would have to be supported by the approved risk-blind
F2-FEHPO design.
The earlier F1-ORD and F2-ADV calculations are historical substrate-specific
contrasts only; neither is a second positive family or can rescue F2-FEHPO.

Two conditions must both hold before useful personalization is claimed:
action-ranking changes replicate across prespecified strata under
\cref{prop:certification}'s procedure, and \cref{eq:personalization} is positive
on the final holdout.  Exploratory heterogeneity may generate later hypotheses
but cannot substitute for either.

\paragraph{Falsification.}  If the registered one-sided test rejects
\eqref{eq:null}, risk ranking is shown to misallocate at at least one registered
capacity.  If it does not reject, ordinal insufficiency has not been established
at the achieved precision; this is not evidence that the orderings are
equivalent.  The insufficiency framing is then withdrawn, with simultaneous intervals
and power reported.  Approximate sufficiency would require a separately
registered equivalence margin and upper-bound procedure.
